% NOETHER PAPER 6 — KOREAN TRANSLATION-PRODUCER DRAFT — UNCHECKED
% Unit: T05-U40; current-authority whole lines 4926--4936
% Source slice: 624 LF UTF-8 bytes; SHA-256 81391C370060148C1F8CC7EC655029726A5D5495F2F68EB3850D97DD7E4B2D92
% Pointer: NOETH-DE-AUTH-v007-20260804; SHA-256 A6A8FC8E5AC24ACAF49DFD55B4B58FA3DA882EF8C3FDD4D136220C8751045156
% German authority: NOETH-DE-ED-0001; 2,153,565 raw bytes; SHA-256 D1F06B311F6CBD991DD247D745DD9A72DDE326A20396DF43CFE0C8EDB1593CDB
% Producer translation only: no source/Korean/formula review, compilation, or rendering.
% Sense window/alternatives: Abbildungssystem→사상계 denotes the §3 lower-variable image system; 대응계 is held for checker comparison. Übertragungsprinzip→전이 원리 keeps the relation-transfer sense, not an operator sense.

그러므로 사상계 $\Ssys^*_{\rho\rho}$에
\[
g_1^*(x),\ g_2^*(x),\ldots,g_k^*(x)
\]
로 주어진 유리기저가 있고, 이것이 관계
\[
f^*(x)=\gamma\bigl(g_1^*(x)\cdots g_k^*(x)\bigr)
\]
를 귀결한다면, 정리 I에 따라 $\Ssys_{n\rho}$에서 관계 (1)이 성립한다. 즉 \emph{유리기저에 대한 전이 원리}는 다음과 같다:

\emph{사상계 $\Ssys^*_{\rho\rho}$에 유리기저가 존재하면 원래 계 $\Ssys_{n\rho}$에도 유리기저가 존재한다.}\srcfn{*)}{여기서 다루는 각 기저 문제에 대해서도 전이 원리는 이에 상응하게 진술된다.}
